\lysh{38791}{4}

\begin{alist}
\item Ο τύπος που εκφράζει τα συνολικά έσοδα $E(x)$ της επιχείρησης από την πώληση $x$ τραπεζιών είναι: $E(x) = p(x) \cdot x = (120 - x) \cdot x = 120x - x^2$.

\item Στον τύπο $K(x) = 20x + 475$, το $475$ δηλώνει ότι υπάρχει ένα σταθερό κόστος $475$ ευρώ, για πάγια έξοδα, όπως π.χ. για το ενοίκιο, το ρεύμα, ενώ το $20$ δηλώνει το κόστος παραγωγής κάθε τραπεζιού (π.χ. το κόστος των υλικών).

\item
\begin{rlist}
\item Τα σημεία $A(5, 575)$ και $B(95, 2375)$ είναι τα δύο σημεία τομής των συναρτήσεων των εσόδων και του κόστους λειτουργίας, δηλαδή τα σημεία στα οποία τα έσοδα της επιχείρησης είναι τα ίδια με το κόστος παραγωγής.
Το σημείο $A(5, 575)$ δείχνει ότι τα έσοδα της επιχείρησης από την πώληση $5$ τραπεζιών είναι ίσα με το κόστος λειτουργίας της, δηλαδή $575$ ευρώ.
Το σημείο $B(95, 2375)$ δείχνει ότι τα έσοδα της επιχείρησης από την πώληση $95$ τραπεζιών είναι ίσα με το κόστος λειτουργίας της, δηλαδή $2375$ ευρώ.
Συμπεραίνουμε ότι και στο σημείο $A$ και στο $B$ η επιχείρηση δεν έχει ούτε κάποιο κέρδος αλλά ούτε και ζημία.

\item Η επιχείρηση έχει κέρδος, αν πουλήσει περισσότερα από $5$ έως λιγότερα από $95$ τραπέζια, δηλαδή από $6$ έως και $94$. Αντίθετα, θα έχει ζημία αν πουλήσει λιγότερα από $5$ ή περισσότερα από $95$ τραπέζια.

\item Αν πουλήσει $110$ τραπέζια, θα έχει ζημία, αφού τότε η τιμή πώλησης θα είναι $p(x) = 120 - 110 = 10$ ευρώ, δηλαδή λιγότερο από τα $20$ ευρώ, που είναι το κόστος για την παραγωγή του ενός τραπεζιού.
\end{rlist}
\end{alist}

